\documentclass[12pt,a4paper]{article}
\usepackage[a4paper,margin=1in]{geometry}
\usepackage{graphicx}
\usepackage{hyperref}
\usepackage{titlesec}
\usepackage{enumitem}
\usepackage{longtable}
\usepackage{booktabs}
\usepackage{fancyhdr}

%----------------------------
% Title and Formatting
%----------------------------
\titleformat{\section}{\large\bfseries}{\thesection}{1em}{}
\titleformat{\subsection}{\normalsize\bfseries}{\thesubsection}{1em}{}
\setlist{nosep}

\pagestyle{fancy}
\fancyhf{}
\rhead{Technical Design Document}
\lhead{\leftmark}
\cfoot{\thepage}

\title{\textbf{Technical Design Document}\\[0.5em]\large Nimbler Sentences With Reinforcement Learning}

\author{Nimble @ WashU Robotics \\ \texttt{washurobotics@gmail.com}}

\author{
    \\
    Tyler Dinh, Lauren Gayle, Anling Zhu, Tenzin Dukdak, \\
    Billy Zhao, Noah Wolk, Mark Seidel, Horlasy Demakpor \\[0.3em]
    Software Team \\[0.2em]
    \texttt{\{dinh.t, gayle.l, anling, dtenzin, billy.z,} \\
    \texttt{n.j.wolk, m.m.seidel, horlasy.d\}@wustl.edu} \\\\
    Elizabeth Joseph, Alexis Luna, Yuxuan Shen, \\
    Kyle Okui, Gabe Minton \\[0.3em]
    Integration Team \\[0.2em]
    \texttt{\{elizabeth.j, alexis.l, notsure, k.okui, g.m.minton\}@wustl.edu}\\\\
}
\date{\today}

%----------------------------
\begin{document}
\maketitle
\newpage
\tableofcontents
\newpage

%----------------------------
\section{Overview}
%Describe the purpose of this document and its intended audience.  
%Mention any prerequisite knowledge required before reading.  
%Specify whether it targets a technical or non-technical audience.
The purpose of this document is to present the technical design of a system that enables a robotic arm to learn and perform handwriting tasks for basic sentences using reinforcement learning (RL). \\\\
This document is intended for a technical audience with a background in computer science, robotics, or machine learning. Non-technical readers may find high-level summaries useful but should refer to supplementary documentation at the end of this document for background on the underlying mechanisms employed.
\newpage
%----------------------------
\section{Problem}

Current robot arms used for writing tasks operate with \textbf{static, pre-programmed motions}. Each specific writing task---whether it's writing a particular sentence, drawing a shape, or creating a signature---requires its own custom motion program. This approach has several critical limitations:

\subsection{The Core Problem}

When a robot needs to write something new, engineers must:
\begin{enumerate}
    \item Manually program the exact trajectory for each letter or word.
    \item Hard-code the sequence of movements.
    \item Test, refine, and repeat this entire process for every new writing task.
\end{enumerate}

A \textbf{dynamic robot arm} that can generalize its writing capabilities would:
\begin{itemize}
    \item \textbf{Scale efficiently}: Write any arbitrary text without reprogramming.
    \item \textbf{Adapt to variations}: Handle different sizes and styles of text.
    \item \textbf{Learn from examples}: Improve through demonstration rather than explicit programming.
\end{itemize}

\subsection{The Core Goal}

The goal is to develop a robot writing system that can:
\begin{itemize}
    \item Accept arbitrary text input (like ``Hello World'' or even more arbitrary, like ''Write me a haiku.'')
    \item Dynamically generate appropriate writing motions.
    \item Execute smooth, legible handwriting without pre-programmed trajectories for each specific string.
\end{itemize}
\newpage
%----------------------------
\section{Tenets}
TODO!
List the key principles or beliefs guiding the design decisions.  
Tenets help align teams and establish common ground on critical design questions.
talk to team about this
\newpage
%----------------------------
\section{Requirements}
TODO! Describe all requirements imposed by the problem.  
Write from the end-user’s perspective. Include user stories or use cases.

\subsection*{Example Use Cases}
\begin{itemize}
  \item As a x, I want to y.
  \item As a x, I want y.
\end{itemize}
%----------------------------
\subsection{Out of Scope}
TODO! List what is explicitly out of scope to avoid misunderstandings.

%----------------------------
\subsection{Success Criteria}
TODO! Describe how success will be measured once the solution is in production.  
Include measurable metrics such as performance, scalability, or cost efficiency.

\subsection*{Example Metrics}
\begin{itemize}
  \item Time-to-market reduced by 90\%.
  \item Supports 10,000 users with 100ms P90 latency.
  \item It can write sentences in English legibly with X\% accuracy.
\end{itemize}
\newpage
%----------------------------
\section{Architecture}
Describe the architecture of the proposed solution using text, diagrams, and bullet points.  
Explain why this design was chosen.

25-26 Architecture

VLA: NORA 1.5
Data Collection Framework: LeRobot
Arm: WidowX AI Trossen Base
CV: OpenCV + Logitech
Additional VLM/LLM for Data Preprocessing: Local Qwen3-4B instance,
Programming Languages: Mainly Python, JSON for data.
Compute: RoboServer

\subsection{Extension of Writing Range with a Linear Actuator}

Writing text is difficult, as the phrase `hello world this is a cool sentence` with the WidowX arm can only be written in the range of action given by the robot and its camera. A linear actuator provides a `sliding window' over the writing surface, enabling arbitrarily long text output.

\subsubsection*{Text Planning}

Before any motion occurs, the full spatial layout of the text should be pre-computed:
\begin{itemize}
  \item Determine character width and inter-character spacing in physical units (cm).
  \item Map the full string to an absolute coordinate sequence along the writing axis.
  \item Determine the total horizontal distance required.
\end{itemize}

\subsubsection*{Text Chunking by Reachability}

Given the current position of the arm, a \textbf{range of action} is defined, which is the zone the writing tool can write in at the current position. The input text to the start of the program is greedily chunked, meaning that characters are grouped until the next character would fall outside the window. Chunk boundaries are placed \emph{between characters}, but should never be chunked mid-character.

\subsubsection*{Slide by Inference on a Chunk}

For each chunk, the VLA (NORA 1.5) executes the writing motions normally. When the chunk is complete, the actuator must activate, moving X distance based on where the next chunk is. Actuator movement is intentionally seperate from VLA inference, as it is purely physical and does not need to be learned. After sliding, a brief re-anchoring step determines where the last written character ended to determine where the arm should resume movement. The VLA is then re-prompted with the next chunk of text.

\subsection{High-Level Overview (HLD)}
List all logical system components:
\begin{itemize}
  \item comp1
  \item comp2
  \item comp3
  \item comp4
  \item comp5
\end{itemize}

Include diagrams as needed:
\begin{center}
%\includegraphics[width=0.8\textwidth]{architecture-diagram.png}
\end{center}

\subsection{API Design}
List all APIs used for interaction between users or services.  
Include HTTP methods, payloads, versions, and example requests/responses.  
Discuss future evolution of the APIs.

\subsection{Data Storage and Model}
Describe the data model and database choice.  
Estimate data volume, forecast growth, and justify scalability.  
Discuss data pipelines, ingestion, and pre-processing if applicable.

\subsection{Application / Component Level Design (LLD)}
Provide detailed design of each system component, including data flow and control flow diagrams.

USE DRAW.IO for this (can install on your computer)

\newpage
%----------------------------
\section{Dependencies}
List external and internal dependencies.  
Document assumptions and risks associated with these dependencies.
%----------------------------
\subsection{Design Alternatives Considered}
Discuss alternative designs evaluated.  
Provide a comparison table highlighting trade-offs.

\begin{longtable}{@{}p{3cm}p{5cm}p{5cm}@{}}
\toprule
\textbf{Option} & \textbf{Pros} & \textbf{Cons} \\ \midrule
Option A & Scalable, cost-efficient & Complex to maintain \\
Option B & Easy to implement & Limited scalability \\ \bottomrule
\end{longtable}

\newpage
%----------------------------
\section{Cost Analysis}
Estimate infrastructure and operational costs.  
Plan for future growth and scalability.

\newpage
%----------------------------
\section{Failures and Risks}
Discuss potential system failures such as:
\begin{itemize}
  \item Dependency failures
  \item Traffic overflow
  \item Performance degradation
  \item Logic bugs
\end{itemize}

Identify risks such as external dependencies or resource limitations.  
Include potential mitigation strategies.

\newpage
%----------------------------
\section{Non-Functional Requirements}
Cover scalability, availability, maintainability, reliability, latency, and security.

\subsection{Scalability}
Expected users, transactions, and data volume.

\subsection{Latency and Availability}
Define SLAs (e.g., P99, P50 latency targets).

\subsection{Maintainability}
Explain maintenance expectations and ownership.

\subsection{Security}
Describe security measures and required compliance levels.

\newpage
%----------------------------
\section{Testing and Observability}
Outline strategies for testing, monitoring, and alerting.

\subsection{Testing}
Explain testing approach: unit, integration, A/B, or stress tests.

\subsection{Metrics and Alarms}
List key performance metrics, dashboards, and alert configurations.

\newpage
%----------------------------
\section{Future Improvements}
List planned features or improvements for future releases.

\newpage
%----------------------------
\section{FAQs}
Include frequently asked questions and their answers.

question?
answer here

\newpage
%----------------------------
\appendix
\section{Appendix A: Subtitle (Optional)}
Include detailed technical analysis or supporting data.

\newpage
%----------------------------
\section{Glossary}
\begin{description}
  \item[ex1] example definition
\end{description}

\newpage
%----------------------------
\section{References}
\begin{itemize}
  \item Service 1: \url{https://example.com}
\end{itemize}

\end{document}
